\section{GPU decode 端到端：常驻权重、KV 布局、sampler 与同步成本}

\subsection{GPU 快，不代表端到端一定快}

GPU 在大 GEMM 上很强，但 decode 阶段常常是小 batch、单 token、长上下文、频繁同步。若只看单个 kernel 时间，很容易误判。端到端 GPU decode 要把权重常驻、KV cache、logits、sampler、copy、stream、event 和 fallback 一起算。

\begin{keyidea}
GPU decode 的核心问题不是“能不能 launch kernel”，而是让权重、KV、activation 和 sampler 尽量留在设备侧，同时把每步同步和 host-device copy 控制到可解释范围。
\end{keyidea}

一个最小 GPU decode step 包括：

\begin{lstlisting}[numbers=none,caption={GPU decode step 账本}]
input token on host
  -> copy token or embedding input to device
  -> run layer kernels
  -> read/write device KV cache
  -> compute logits on device
  -> sampler on device or copy logits to host
  -> emit token event on host
\end{lstlisting}

每个箭头都有成本。若每步把所有中间 activation 拷回 CPU 做调试，GPU kernel 再快也没意义。debug oracle 可以按 trace level 打开，但默认路径应尽量减少同步。

\subsection{权重常驻设备}

GPU 后端的基本要求是权重 prepare 后常驻设备。每次请求或每个 token 重新上传权重，端到端必然崩溃。

\begin{lstlisting}[numbers=none,caption={GPU prepare 阶段}]
load model files on host
verify manifest and quant descriptors
pack or convert weights for GPU backend
allocate DeviceBuffer for each weight group
copy weights to device
record device_weight_table in plan report
\end{lstlisting}

plan report 要记录 device weight 字节、上传时间、quant format 和 fallback。若 GPU 显存不够，应该在 compile/prepare 阶段拒绝或选择 CPU fallback，而不是在请求运行中途失败。

\subsection{activation 和 workspace}

decode 每层会产生临时 activation。GPU 后端应使用 workspace 或 arena，避免每个 step 调用设备分配器。

\begin{lstlisting}[numbers=none,caption={GPU workspace 字段}]
GpuWorkspace:
  activation_buffers
  attention_scratch
  logits_buffer
  sampler_temp
  size_bytes
  alignment
  stream_owner
\end{lstlisting}

workspace 的生命周期通常是 plan 或 session。若 workspace 被多个 session 共享，调度器必须保证不会并发写同一 buffer。若每个 session 独占 workspace，并发容量要把它算进 admission control。

\subsection{Device KV cache}

若 layer kernel 在 GPU 上运行，KV cache 最好也在 GPU 上。否则每层 attention 都要在主机和设备之间传 K/V，成本不可接受。

\begin{lstlisting}[numbers=none,caption={GPU KV cache 设计点}]
DeviceKvCache:
  layout: contiguous | paged
  dtype: fp16 | int8 | int4
  block_table_optional
  scale_table_optional
  max_context
  resident_device
\end{lstlisting}

连续 KV 在 GPU 上地址简单，但并发浪费显存；分页 KV 更灵活，但 kernel 要通过 block table 找 K/V。长上下文下，attention kernel 的读取模式决定了 coalescing。若每个 thread 读取不连续地址，带宽会下降。

\subsection{GPU executable plan}

GPU 后端也需要 executable plan，而不是在 runtime 中临时拼 kernel。

\begin{lstlisting}[numbers=none,caption={GPU segment descriptor}]
GpuSegmentDesc:
  segment_id
  op_or_fused_region
  kernel_name
  grid_shape
  block_shape
  input_buffers
  output_buffers
  state_reads
  state_writes
  workspace_range
  stream_policy
  fallback_segment_optional
\end{lstlisting}

\code{state\_reads} 和 \code{state\_writes} 对 attention 很重要：decode attention 读历史 KV，同时写当前 K/V。若 plan 不表达状态边，调度器可能错误重排 segment；若 workspace range 不明确，多个 segment 可能覆盖同一临时区域；若 fallback segment 不记录，性能报告会失真。

\subsection{device residency report}

GPU prepare 完成后，要输出哪些对象在设备上。

\begin{lstlisting}[numbers=none,caption={DeviceResidencyReport}]
DeviceResidencyReport:
  model_id
  device_id
  weights:
    total_bytes
    quant_profile
    packed_layout
  kv_cache:
    layout
    max_reserved_bytes
    dtype
  workspace:
    bytes
  logits:
    device_buffer_bytes
  fallback:
    unsupported_segments
\end{lstlisting}

这份报告能提前发现“显存看似够，实际 packed weight 加 workspace 不够”的问题。它也能解释为什么同一模型在不同 GPU 上 plan 不同：某些设备支持某种低比特 kernel，另一些设备只能 fallback。

\subsection{prefill 和 decode 的 GPU 策略不同}

prefill 有较大的 token 维度，适合 GEMM 和 batched attention。decode 单 token 时，kernel launch、KV 读流和小矩阵效率更突出。

\begin{center}
\begin{tabular}{p{0.18\linewidth}p{0.34\linewidth}p{0.32\linewidth}}
\toprule
阶段 & GPU 机会 & GPU 风险 \\
\midrule
prefill & 大 GEMM、并行 attention、吞吐高 & 显存峰值、长 prompt attention \\
decode & 多 session 合批、KV 常驻、低同步 & 小 batch 低占用、launch/copy 成本 \\
\bottomrule
\end{tabular}
\end{center}

因此 GPU 报告必须拆 prefill 和 decode。只报总 tokens/s 会掩盖问题：一个后端可能 prefill 非常快，但单用户 decode 受 launch 成本限制；也可能多 session 合批后吞吐高，但 p95 变差。

\subsection{sampler 放设备侧的收益和代价}

若 logits 在 GPU 上，CPU sampler 需要把 logits 或 top-k 候选拷回 host。完整词表 logits 可能是几十 KB 到几百 KB，每步还会同步 stream。设备侧 sampler 可以减少同步，但要实现 top-k/top-p、随机数和 stop 检查。

折中做法是 GPU 先选 top-k 候选，只拷回候选 id 和分数给 CPU 最终采样。这样减少拷贝，但仍保留业务逻辑和可复现随机数在 CPU。

\begin{center}
\begin{tabular}{p{0.22\linewidth}p{0.34\linewidth}p{0.28\linewidth}}
\toprule
策略 & 拷贝量 & 实现复杂度 \\
\midrule
拷完整 logits 到 CPU & 高 & 低 \\
GPU top-k，CPU 采样 & 中低 & 中 \\
全 GPU sampler & 低 & 高 \\
\bottomrule
\end{tabular}
\end{center}

第一版可以先用 CPU sampler，但必须在 trace 中单独记录 \code{logits\_copy\_ns} 和 \code{sampler\_ns}。这样把 sampler 下沉到 GPU 时，收益能被证明。

\subsection{stream 和 event}

GPU 后端不能只用一个全局同步。最小资源模型：

\begin{lstlisting}[numbers=none,caption={GPU 执行资源}]
GpuExecutionContext:
  device_id
  stream
  events:
    prefill_start
    prefill_end
    decode_step_start
    decode_step_end
  workspace
  error_state
\end{lstlisting}

event 用来测量 GPU 时间，host 计时用来测端到端时间。二者都要记录。若 GPU event 显示 kernel 很快，但 host 端 decode step 慢，问题可能在 copy、sync、sampler、队列或事件等待。

\subsection{同步点要显式列出}

GPU 性能问题经常藏在同步点。常见同步来源：

\begin{itemize}
  \item 拷回 logits 给 CPU sampler。
  \item debug 模式 materialize 中间张量。
  \item host 等待 GPU event 才能发 token。
  \item fallback segment 在 CPU 和 GPU 之间来回搬数据。
  \item 多 stream 之间的依赖 event。
  \item device allocator 或临时 buffer 初始化。
\end{itemize}

每个同步点都要进入 trace。若只是看 kernel profile，可能看不到 host 在等待。端到端优化时，减少一个同步点往往比把某个小 kernel 优化 10\% 更有效。

\subsection{合批 decode}

GPU decode 要提高吞吐，通常需要合批多个 session 的当前 token。合批后，一个 step 的输入形状从 \code{[1,H]} 变成 \code{[B,H]}，权重 GEMV 变成小 GEMM，GPU 利用率提高。但合批有代价：请求要等待同批其他 session，KV cache 来自不同 session，stop 和 sampler 也各自独立。

\begin{lstlisting}[numbers=none,caption={GPU decode 合批元数据}]
DecodeBatch:
  batch_size
  session_ids
  model_positions
  kv_block_tables
  input_token_ids
  sampler_states
\end{lstlisting}

合批不是把 token 拼起来就完事。每个 session 的 \code{model\_position}、KV block table、stop 条件和 sampler state 都不同。kernel 可以共享权重，但不能混淆状态。调度器需要设置最大等待时间，否则为了凑 batch 会拉高 p95。

\subsection{GPU oracle 策略}

GPU 路径也要有 oracle，但不能每步都拷完整中间张量。可用策略：

\begin{itemize}
  \item smoke 模式：只比较最终 logits 和 top-k。
  \item checkpoint 模式：抽样比较某几层输出。
  \item debug 模式：materialize 指定 segment 的输入输出。
  \item byte 模式：只比较 packed weight 或量化解包。
\end{itemize}

trace level 控制 oracle 成本。默认生产路径不应频繁同步中间张量；调试时可以打开指定 segment 的 checkpoint。这样既保留定位能力，又不让 oracle 扭曲性能数据。

\subsection{fallback 不是静默降级}

GPU 可能不支持某个 dtype、某个 group size、某个上下文长度或某个算子融合。fallback 可以是好策略，但必须显式记录：

\begin{lstlisting}[numbers=none,caption={GpuFallbackEvent 字段}]
GpuFallbackEvent:
  segment_id
  requested_backend
  actual_backend
  reason:
    dtype_unsupported
    group_size_unsupported
    workspace_exceeded
    kernel_missing
    device_memory_exceeded
  performance_impact_estimate
\end{lstlisting}

静默 fallback 会制造严重误导：用户以为在测 GPU，实际跑 CPU；或者某个 segment 来回拷贝，端到端比纯 CPU 更慢。

\subsection{GPU StepTrace}

GPU decode 每步至少要记录：

\begin{lstlisting}[numbers=none,caption={GPU StepTrace 字段}]
GpuStepTrace:
  request_id
  step
  context_length
  batch_size
  stream_id
  kernel_time_ns
  copy_h2d_bytes
  copy_d2h_bytes
  copy_time_ns
  sync_time_ns
  sampler_time_ns
  fallback_count
  device_memory_peak
\end{lstlisting}

这些字段能定位典型问题：kernel 时间低但同步高，说明 host-device 边界有问题；copy 字节突然上升，说明某个 debug 或 fallback materialize 了中间张量；batch size 太小，说明调度器没有把 decode 合并起来。

\subsection{为什么 decode 小 batch 难吃满 GPU}

GPU 适合大量独立工作并行推进。prefill 有较大的 token 维度，linear 和 attention 可以形成大 GEMM 或大 batch attention；decode 单用户、单 token 时，很多算子形状退化成 \code{[1,H]} 乘权重，或者一个 query 读长 KV。问题不是 GPU 算不动，而是并行度、数据复用和固定开销不匹配。

可以把 GPU decode 的一个 step 拆成四类成本：

\begin{center}
\begin{tabular}{p{0.22\linewidth}p{0.36\linewidth}p{0.28\linewidth}}
\toprule
成本 & 来源 & 小 batch 风险 \\
\midrule
kernel launch & host 提交 kernel 到 device & 每步许多小 kernel，固定开销占比高 \\
global memory & 权重、KV、activation、logits 读写 & 权重/KV 字节大，算术强度低 \\
occupancy & SM 上活跃 warp 数和资源使用 & grid 太小或寄存器/shared memory 占用高 \\
sync/copy & logits、sampler、fallback、debug materialize & host 等待 device，stream 被切断 \\
\bottomrule
\end{tabular}
\end{center}

因此“GPU 峰值 TFLOP/s 很高”不能直接推出“单用户 decode 很快”。报告必须拆 prefill 和 decode，还要记录 batch size、context length、kernel 数量、launch/sync/copy 时间和 device 端实际 kernel 时间。

\subsection{occupancy 的原理边界}

\term{occupancy}{占用率} 粗略表示一个 SM 上同时驻留的 active warps 与硬件上限的比例。它受 block size、每线程寄存器数量、每 block shared memory、线程块数量和硬件限制影响。更高 occupancy 能帮助隐藏 memory latency，但不是越高越快；如果为了提高 occupancy 减少寄存器导致 spill，或让每个线程做更少有用工作，性能可能下降。

decode 小 batch 常见的 occupancy 问题有两类。第一，grid 本身太小：batch=1，某些小算子只发出很少 block，很多 SM 没有工作。第二，单个 kernel 资源占用高：attention kernel 为了处理长 context 使用较多寄存器或 shared memory，每个 SM 能驻留的 block 少，隐藏延迟能力下降。

\begin{lstlisting}[numbers=none,caption={occupancy 报告要保存的字段}]
kernel_name
grid_dim
block_dim
registers_per_thread
shared_memory_per_block
estimated_active_warps_per_sm
achieved_occupancy_if_available
dram_bytes
elapsed_us
\end{lstlisting}

occupancy 只能解释“是否有足够并行工作隐藏等待”，不能单独证明瓶颈。若 achieved occupancy 很高但 DRAM bandwidth 接近上限，继续提高 occupancy 不会有太大收益；若 occupancy 低且每步 kernel 很短，可能要做算子融合、合批或 persistent kernel，而不是只调 block size。

\subsection{tensor core 为什么更偏爱 prefill}

Tensor core 擅长固定 tile 的矩阵乘加，例如 FP16/BF16/TF32/INT8 的矩阵 tile。prefill 的 \code{A[T,H] x W[O,H]} 中，\code{T} 足够大时，可以形成大矩阵 tile，权重和 activation 都能在 tile 内复用。decode batch=1 时，\code{A[1,H]} 让矩阵的一维太小，tile 利用率低；即使使用 tensor core，很多 lane 可能没有足够有效工作。

\begin{center}
\begin{tabular}{p{0.20\linewidth}p{0.32\linewidth}p{0.32\linewidth}}
\toprule
阶段 & tensor core 机会 & 常见限制 \\
\midrule
prefill & 大 GEMM，tile 饱满，权重复用高 & activation 峰值和 attention workspace \\
decode 单用户 & 小 GEMM/GEMV，tile 不饱满 & launch、权重流、KV 读、低 batch \\
decode 合批 & \code{B x H} 变大，可转小 GEMM & 等待凑 batch，p95/SLO 压力 \\
\bottomrule
\end{tabular}
\end{center}

所以 GPU serving 的关键是调度：在不破坏 SLO 的前提下，把多个 session 的 decode step 合并成更适合 GPU 的小 batch。batch 大了，吞吐可能上升；等待时间也可能上升。报告必须同时写 tokens/s、TTFT、TPOT、p50/p95/p99，而不是只写平均吞吐。

\subsection{memory bandwidth 下界}

延续 CPU 章节的 7B KV 账本。若 context=4096，fp16 KV 对一个 session 约 512 MiB。GPU attention kernel 不一定每 step 从 DRAM 读完整 512 MiB；cache、分块和并行归约会改变实际流量。但下界分析仍然有用：只要某个 step 的有效 KV 读取接近数百 MiB，latency 就会受到显存带宽约束。

假设有效读流为 \code{400 MiB}，实际可用带宽 \code{800 GB/s}：

\[
time_{min} \approx 0.4 / 800 = 0.0005s = 0.5ms
\]

这只是 KV 读取，不含 projection、MLP、lm head、softmax、kernel launch 和 sampler。如果 profiler 显示 kernel 时间低于这个数量级，需要核对实际读字节、是否使用了更低精度 KV、是否只读窗口、是否 batch 共享了某些 prefix、计时边界是否只覆盖了部分 kernel。

\subsection{kernel launch overhead 和融合}

decode 每个 token 要经过多层，每层又有 RMSNorm、QKV projection、RoPE、attention、output projection、MLP、residual 等操作。如果每个小操作都是独立 kernel，launch overhead 和 stream 同步会变成显著成本。融合的目标不是让代码更炫，而是减少中间内存流和 kernel 数量。

\begin{center}
\begin{tabular}{p{0.25\linewidth}p{0.33\linewidth}p{0.27\linewidth}}
\toprule
融合方向 & 可能收益 & 风险 \\
\midrule
RMSNorm + QKV 前处理 & 减少 activation 读写 & kernel 专用化增加 \\
RoPE + Q/K 写 cache & 减少一次 K/Q materialize & position/layout bug 更难查 \\
attention score + softmax + V reduce & 减少中间 score 存储 & 数值稳定和 debug 成本 \\
bias + activation + residual & 减少小 kernel 和内存流 & tail/广播规则复杂 \\
\bottomrule
\end{tabular}
\end{center}

融合前必须有 reference 和 checkpoint oracle。否则 fused kernel 输出错了，很难判断是 RoPE position、softmax 稳定性、KV address、tail mask 还是 dtype 转换出错。

\subsection{小 batch GPU decode 的实验矩阵}

GPU decode 报告至少要做 shape sweep，而不是只测一个 prompt：

\begin{center}
\begin{tabular}{p{0.22\linewidth}p{0.30\linewidth}p{0.32\linewidth}}
\toprule
变量 & 建议取值 & 观察目的 \\
\midrule
batch size & 1, 2, 4, 8, 16 & 找到吞吐和 p95 的拐点 \\
context length & 128, 512, 2048, 4096 & 分离 KV 读流影响 \\
kernel fusion & off/on & 判断 launch 和中间内存流成本 \\
sampler 位置 & CPU, GPU top-k, GPU sampler & 判断 copy/sync 成本 \\
KV dtype & fp16, int8 if available & 判断容量、带宽和误差权衡 \\
\bottomrule
\end{tabular}
\end{center}

每组都要输出 \code{GpuStepTrace}，至少包含 kernel time、copy time、sync time、sampler time、fallback count、batch wait time 和 p95/p99。只有这样才能回答：GPU 慢是 kernel 本身慢，还是 batch 太小、launch 太多、copy 太多、fallback 静默发生，或者 scheduler 为吞吐牺牲了尾延迟。

\subsection{GPU 原理展开模板}

本书只把 GPU 作为 AI Infra 后端边界，不把它写成完整 CUDA 教材。但凡本书提到 GPU kernel，都必须按下面模板说明：

\begin{itemize}
  \item thread/block/warp/SM 如何映射到 tensor 维度。
  \item global memory 读写字节和 coalescing 条件。
  \item shared memory 或寄存器保存了什么 tile。
  \item occupancy 受哪些资源限制。
  \item 是否使用 tensor core，tile 是否饱满。
  \item kernel launch、copy、sync 是否进入端到端 trace。
  \item 小 batch、长 context 和 fallback 的边界是什么。
\end{itemize}

\subsection{发布前 GPU 检查表}

\begin{rubricbox}
\begin{itemize}
  \item prepare 阶段完成权重上传，decode step 不重复上传权重。
  \item device residency report 记录权重、KV、workspace、logits 和 fallback。
  \item GPU segment descriptor 记录 kernel、buffer、state read/write、workspace 和 stream。
  \item StepTrace 同时有 GPU event 时间和 host 端 elapsed 时间。
  \item logits copy、sampler、sync 和 fallback 单独计时。
  \item 合批 decode 记录 batch size、等待时间和每个 session 的 position。
  \item debug oracle 不污染默认性能报告。
\end{itemize}
\end{rubricbox}

\subsection{本节验收线}

读完本节，应该能设计 GPU decode 的端到端合同：

\begin{itemize}
  \item 权重在 prepare 后常驻 device，并写入 plan report。
  \item activation 和 workspace 不在每 step 动态分配。
  \item KV cache 在 device 侧维护，连续或分页布局都有明确地址合同。
  \item prefill 和 decode 分开报告，不用一个 tokens/s 混过。
  \item sampler、logits copy、stream sync 和 fallback 都进入 StepTrace。
  \item fallback 必须显式记录原因和实际 backend。
  \item 用 GPU event 和 host trace 同时解释 kernel 时间与端到端时间。
\end{itemize}
